\chapter{Related Work}\label{ch:related}

% In this chapter, you describe how other people have dealt with similar problems that you are dealing with.
% What were their techniques?
% What were their results?
% What observations did they make that you have used?
% In what way is your work different?

% The role of the related work is threefold:
% \begin{enumerate}
% \item Your work is grounded in earlier work done by other researchers.
% That means you use techniques and methods, ways of measuring and evaluating, datasets and terminology used by others working on the same or a related problem before you.
% \item You want to compare your results with those of others.
% This is especially true for key papers with respect to your research, \ie{}~\ac{SOTA}.
% Those papers need to be referred to in this chapter.
% This could also lead to a hypothesis, which is the answer that you expect based on previous literature.
% \item The previous two points mentioned cover the main function of this chapter.
% In addition with related work, you inform the reader on the \ac{SOTA}, the state of the problem, and the
% existing solutions.
% You also show that you know what you are talking about, that you know what is ``for sale'', and that you use the best there is.
% \end{enumerate}

% Depending on your story, you may want to put this chapter at the end of your thesis (before the Conclusions chapter) or at the beginning of your thesis (after the Introduction chapter).
% Typically, when your work is a direct extension or modification of existing work, you want to put this at the beginning to compare and contrast your research.
% This helps in showing the reader the context and demonstrating the novelty of your approach.

% When your research is in a very new area where you explore a novel concept, you may want to put the Related Work at the end.
% It is beneficial to first get the reader comfortable with your research before showing how it fits within the wider research area.
% Showing the related work first may distract the reader and make them lose interest in the rest of the thesis.

% In the end, think about the story you want to tell and discuss this with your academic supervisor.

This chapter positions the thesis with respect to prior work. The relevant literature spans four areas: explainable artificial intelligence, software architecture and abstraction, usable security and security communication, and natural language generation or data-to-text. Each area contributes part of the foundation needed for this thesis, but none of them, on its own, fully addresses the problem studied here. The main argument developed in this chapter is that current work remains fragmented: explainability research is often model-centred, cybersecurity tooling is largely technical-user oriented, and software architecture research rarely treats explainability for non-technical stakeholders as a first-class design problem. That gap motivates the contribution of this thesis.

\section{Explainable AI and Human-Centred Explanation}

The recent wave of XAI research has produced a broad range of methods for making model behaviour more understandable. Surveys typically distinguish among feature attribution, surrogate models, example-based explanations, and counterfactual approaches \cite{guidotti2018survey, barredo2020xai}. Widely cited methods such as LIME and SHAP have become common reference points because they offer general-purpose approaches for explaining black-box predictions \cite{ribeiro2016should, lundberg2017unified}. Counterfactual explanations have further expanded the space by showing how inputs would need to change to obtain a different prediction \cite{wachter2017counterfactual}. These methods have been influential because they provide practical tools for opening up model outputs that would otherwise remain opaque.

At the same time, the XAI literature has also made clear that explanation quality cannot be judged independently of the user and task. Doshi-Velez and Kim argue that interpretability should be evaluated in relation to its intended use \cite{doshivelez2017rigorous}. Lipton similarly criticises the tendency to treat interpretability as a single, self-evident property and shows that different explanation goals are often conflated under the same label \cite{lipton2018mythos}. Miller goes further by drawing on social science research to show that human explanations are selective and contrastive, meaning that people usually want explanations that answer a specific practical question rather than an exhaustive account of internal model mechanics \cite{miller2019explanation}.

These observations are highly relevant to this thesis, but they also point to a limitation in the current state of the art. Most XAI research remains focused on explaining the output of a machine learning model in isolation. The typical unit of analysis is a prediction. This is useful when the problem is to understand a classifier, but it is less suitable when the practical problem is to design a broader software workflow in which outputs from one or more tools must be translated into stakeholder-oriented representations. In vulnerability management, the challenge is not simply that a model is opaque. The challenge is that technical findings need to be transformed into explanations that support communication, prioritisation, and decision-making across roles. Existing XAI methods may contribute ingredients to such a workflow, but they do not by themselves provide the architectural guidance needed to integrate explanation into the system as a maintainable and traceable responsibility.

A second limitation is that XAI methods often presuppose a technically informed user. Even when explanations are labelled as interpretable, they may still require familiarity with machine learning concepts, feature spaces, or model behaviour. For non-technical stakeholders, an attribution score or local surrogate may still be too far removed from the practical question they need answered. This is especially problematic in cybersecurity, where the decision context often includes operational urgency, uncertainty, and the need to justify actions to people outside the security function.

\section{Software Architecture, Abstraction, and Quality Attributes}

Software architecture research offers concepts that are directly relevant to this thesis, even though it rarely addresses explainability in the specific sense studied here. Architecture has long been concerned with managing complexity, separating concerns, and providing structure that supports evolution and communication \cite{bass2012software}. Parnas' classic work on modularisation established that effective decomposition should be driven by information hiding and design decisions, not merely by task breakdown \cite{parnas1972modules}. Kruchten’s 4+1 view model further formalised the idea that systems should be represented through multiple views because different stakeholders need different forms of understanding \cite{kruchten1995viewmodel}.

These contributions are important because they provide the conceptual basis for treating explanation as an architectural concern. If different stakeholders need different views of the same underlying system, then cybersecurity findings should not necessarily be represented in one uniform form. Likewise, if information hiding and separation of concerns improve maintainability, then explanation logic should not be embedded in ad hoc reports or informal analyst interpretation. Instead, explanation generation, abstraction, traceability, and presentation can be modelled as explicit architectural responsibilities.

However, the architecture literature generally stops short of this application. It gives strong guidance on how to structure systems for qualities such as modifiability, usability, performance, and maintainability \cite{bass2012software}, but relatively little work applies these principles to the design of stakeholder-oriented explanation systems in cybersecurity. Architectural views are usually discussed in relation to developers, maintainers, operators, or other technical stakeholders. Less attention has been paid to systems where one of the central architectural problems is how to transform technical outputs into understandable representations for non-technical users.

This is where the current literature becomes insufficient for the thesis problem. The architecture field provides the principles, but not the domain-specific adaptation. It suggests that multiple views and abstraction layers are useful, but it does not tell us how to structure an explanation pipeline for vulnerability reporting, how to preserve traceability to technical evidence, or how to balance explanation richness against maintainability. Those questions remain open.

\section{Usable Security and Security Communication}

Usable security research has shown for many years that security systems cannot be evaluated only in technical terms. If a tool generates information that users cannot correctly interpret, the overall effectiveness of the security process is reduced. Cranor’s work on the human in the loop formalises this broader perspective by showing that security outcomes depend on human understanding, attention, and action, not only on the technical correctness of the mechanism \cite{cranor2008framework}.

Work on intrusion detection systems illustrates this well. Werlinger et al.\ show that even experienced practitioners face considerable difficulty when deploying, configuring, and interpreting IDS outputs \cite{werlinger2008ids}. Their findings are important for this thesis because they shift the focus away from a simplistic assumption that better security tools automatically produce better decisions. In reality, interpretation itself is labour-intensive and socially distributed across roles. Decisions depend on expertise, context, and communication among stakeholders.

This literature strongly supports the thesis motivation, but it also leaves room for contribution. Usable security research is very effective at diagnosing human-facing problems in security systems, yet it often remains descriptive or interface-oriented. It identifies pain points such as cognitive overload, misinterpretation, and configuration burden, but less often proposes a software architectural response to those problems. In other words, it tells us that communication matters, but not in detail how explanation responsibilities should be structured across the system so that outputs become more consistent, role-specific, and maintainable over time.

Cybersecurity dashboards and reporting tools partially address this issue by aggregating and visualising findings, but they often remain close to the original technical indicators. In practice, they improve access to information more than they improve explanation. This matters because stakeholder communication in vulnerability management requires more than a visual summary. It requires a rationale that connects technical evidence to why a risk matters, why it is prioritised, and what should happen next.

\section{Data-to-Text and Translation of Technical Information}

Natural language generation and data-to-text research offer a more directly relevant perspective on the translation problem. Reiter and Dale describe how structured input can be converted into human-readable text through content selection, aggregation, and linguistic realisation \cite{reiter2000nlg}. Gatt and Krahmer survey a mature body of work showing that such methods can support understanding and decision-making in domains where users do not interact directly with the raw underlying data \cite{gatt2018survey}.

This literature is attractive for the current thesis because it reframes explanation as a transformation problem. Instead of exposing all technical details, a system can select what is relevant, group it into higher-level statements, and present it in a form appropriate for the audience. In vulnerability management, this suggests a way to move from raw scanner output to short risk summaries, structured rationales, and actionable recommendations.

At the same time, this body of work has a limitation when brought into cybersecurity. Most data-to-text systems are not designed around traceability to security evidence, role-specific explanation needs, or the quality attributes that become important when explanations are part of a software workflow that must evolve over time. They provide useful design strategies, but not a complete answer to the architectural problem. In other words, they are highly relevant as a source of explanation patterns, but they do not solve the integration challenge that this thesis addresses.

\section{Research Gap}

Taken together, the literature reveals a clear gap.

XAI research provides methods for explaining model outputs, but most of it remains focused on model-level transparency rather than on explanation as a system-level design concern. Software architecture research provides principles for abstraction, modularisation, and stakeholder views, but it rarely applies them to the design of explanation pipelines for cybersecurity communication. Usable security research demonstrates that security outputs are difficult to interpret and that human understanding matters, but it does not usually translate that insight into concrete architectural designs. Data-to-text research shows how structured data can be turned into readable narratives, yet it is not primarily concerned with cybersecurity workflows, traceability requirements, or long-term maintainability of explanation logic.

As a result, existing work does not sufficiently address the following combined problem: how to design a software architecture that transforms technical vulnerability findings into stakeholder-oriented explanations that are understandable, actionable, and traceable to the underlying evidence.

That is the gap this thesis addresses. Rather than proposing yet another standalone explanation method, the thesis contributes a Software Engineering approach that treats explanation as an architectural responsibility. The focus is on designing abstractions, components, and flows that support stakeholder-oriented cybersecurity reporting in a vulnerability management workflow. In this way, the thesis connects areas that are currently treated separately in the literature and provides a more integrated design-oriented response to the problem.